% GENERATED FILE. Edit canonical lessons, then run npm run generate:books.
% canonical-source-hash: fnv1a64:2de2e5d57fb5ce60
% canonical-lessons: ML-C190-pakarttuka, ML-C190-olikkuka2, ML-C190-nekkuka, ML-C190-kulaykkuka, ML-C190-arruka

\chapter{To copy, To hide, To press, To knead, To cool}
\label{ch:ml-to-copy-to-hide-to-press-to-knead-to-cool}
\hlchaptermodality{190}

\begin{chapteropening}
\textbf{By the end of this chapter:} \emph{I can say to copy, to hide, to press, to knead and to cool.}

Complete the last lesson of chapter 190: I can say to copy, to hide, to press, to knead and to cool.
\end{chapteropening}

\section[\texorpdfstring{\ml{പകർത്തുക} (pakarttuka)}{പകർത്തുക (pakarttuka)}]{\ml{പകർത്തുക} (pakarttuka) — to copy}

\label{lesson:ML-C190-pakarttuka}



\begin{quote}
Before the new one: say the Malayalam for to roll, then the Malayalam for to crawl.
\end{quote}

\subsection*{You'll want to know: \ml{പകർത്തുക}}
\textbf{\ml{പകർത്തുക}} — \emph{pakarttuka} — \textquotedblleft{}to copy\textquotedblright{}.

\begin{grammarlens}[title={What you've built}]
One more word for this chapter.
\end{grammarlens}

\subsection*{Your turn}
\begin{itemize}
\raggedright
  \item \emph{Say it:} \emph{pakarttuka}
  \item \emph{Say it:} \emph{pakarttuka}, once more
  \item \emph{Say it:} say \emph{pakarttuka}
  \item \emph{Recall:} say the Malayalam for to roll, then the Malayalam for to crawl, then say \emph{pakarttuka} again
\end{itemize}

\begin{tcolorbox}[breakable,colback=teal!4,colframe=teal!35!black,title={Before you move on}]
What does \ml{പകർത്തുക} mean? (\textquotedblleft{}To copy\textquotedblright{}.) Say it once more.
\end{tcolorbox}
\section[\texorpdfstring{\ml{ഒളിക്കുക} (oḷikkuka)}{ഒളിക്കുക (oḷikkuka)}]{\ml{ഒളിക്കുക} (oḷikkuka) — to hide (oneself)}

\label{lesson:ML-C190-olikkuka2}



\begin{quote}
Before the new one: say the Malayalam for to crawl, then the Malayalam for to copy.
\end{quote}

\subsection*{You'll want to know: \ml{ഒളിക്കുക}}
\textbf{\ml{ഒളിക്കുക}} — \emph{oḷikkuka} — \textquotedblleft{}to hide (oneself)\textquotedblright{}.

\begin{grammarlens}[title={What you've built}]
One more word for this chapter.
\end{grammarlens}

\subsection*{Your turn}
\begin{itemize}
\raggedright
  \item \emph{Say it:} \emph{oḷikkuka}
  \item \emph{Say it:} \emph{oḷikkuka}, once more
  \item \emph{Say it:} say \emph{oḷikkuka}
  \item \emph{Recall:} say the Malayalam for to crawl, then the Malayalam for to copy, then say \emph{oḷikkuka} again
\end{itemize}

\begin{tcolorbox}[breakable,colback=teal!4,colframe=teal!35!black,title={Before you move on}]
What does \ml{ഒളിക്കുക} mean? (\textquotedblleft{}To hide (oneself)\textquotedblright{}.) Say it once more.
\end{tcolorbox}
\section[\texorpdfstring{\ml{ഞെക്കുക} (ñekkuka)}{ഞെക്കുക (ñekkuka)}]{\ml{ഞെക്കുക} (ñekkuka) — to press (a button)}

\label{lesson:ML-C190-nekkuka}



\begin{quote}
Before the new one: say the Malayalam for to copy, then the Malayalam for to hide.
\end{quote}

\subsection*{You'll want to know: \ml{ഞെക്കുക}}
\textbf{\ml{ഞെക്കുക}} — \emph{ñekkuka} — \textquotedblleft{}to press (a button)\textquotedblright{}.

\begin{grammarlens}[title={What you've built}]
One more word for this chapter.
\end{grammarlens}

\subsection*{Your turn}
\begin{itemize}
\raggedright
  \item \emph{Say it:} \emph{ñekkuka}
  \item \emph{Say it:} \emph{ñekkuka}, once more
  \item \emph{Say it:} say \emph{ñekkuka}
  \item \emph{Recall:} say the Malayalam for to copy, then the Malayalam for to hide, then say \emph{ñekkuka} again
\end{itemize}

\begin{tcolorbox}[breakable,colback=teal!4,colframe=teal!35!black,title={Before you move on}]
What does \ml{ഞെക്കുക} mean? (\textquotedblleft{}To press (a button)\textquotedblright{}.) Say it once more.
\end{tcolorbox}
\section[\texorpdfstring{\ml{കുഴയ്ക്കുക} (kuḻaykkuka)}{കുഴയ്ക്കുക (kuḻaykkuka)}]{\ml{കുഴയ്ക്കുക} (kuḻaykkuka) — to knead (dough)}

\label{lesson:ML-C190-kulaykkuka}



\begin{quote}
Before the new one: say the Malayalam for to hide, then the Malayalam for to press.
\end{quote}

\subsection*{You'll want to know: \ml{കുഴയ്ക്കുക}}
\textbf{\ml{കുഴയ്ക്കുക}} — \emph{kuḻaykkuka} — \textquotedblleft{}to knead (dough)\textquotedblright{}.

\begin{grammarlens}[title={What you've built}]
One more word for this chapter.
\end{grammarlens}

\subsection*{Your turn}
\begin{itemize}
\raggedright
  \item \emph{Say it:} \emph{kuḻaykkuka}
  \item \emph{Say it:} \emph{kuḻaykkuka}, once more
  \item \emph{Say it:} say \emph{kuḻaykkuka}
  \item \emph{Recall:} say the Malayalam for to hide, then the Malayalam for to press, then say \emph{kuḻaykkuka} again
\end{itemize}

\begin{tcolorbox}[breakable,colback=teal!4,colframe=teal!35!black,title={Before you move on}]
What does \ml{കുഴയ്ക്കുക} mean? (\textquotedblleft{}To knead (dough)\textquotedblright{}.) Say it once more.
\end{tcolorbox}
\section[\texorpdfstring{\ml{ആറ്റുക} (āṟṟuka)}{ആറ്റുക (āṟṟuka)}]{\ml{ആറ്റുക} (āṟṟuka) — to cool (a hot drink) by pouring}

\label{lesson:ML-C190-arruka}



\begin{quote}
Before the new one: say the Malayalam for to press, then the Malayalam for to knead.
\end{quote}

\subsection*{You'll want to know: \ml{ആറ്റുക}}
\textbf{\ml{ആറ്റുക}} — \emph{āṟṟuka} — \textquotedblleft{}to cool (a hot drink) by pouring\textquotedblright{}.

\begin{grammarlens}[title={What you've built}]
One more word for this chapter.
\end{grammarlens}

\subsection*{Your turn}
\begin{itemize}
\raggedright
  \item \emph{Say it:} \emph{āṟṟuka}
  \item \emph{Say it:} \emph{āṟṟuka}, once more
  \item \emph{Say it:} say \emph{āṟṟuka}
  \item \emph{Recall:} say the Malayalam for to press, then the Malayalam for to knead, then say \emph{āṟṟuka} again
\end{itemize}

\begin{tcolorbox}[breakable,colback=teal!4,colframe=teal!35!black,title={Before you move on}]
What does \ml{ആറ്റുക} mean? (\textquotedblleft{}To cool (a hot drink) by pouring\textquotedblright{}.) Say it once more.
\end{tcolorbox}
